\SourceLeaf{141}%
Sed hic quoque praestat, loco tertiae classis binas stabilire, ut omnino habeantur quaternae.

Assumta radice quacunque primitiva pro basi, residua biquadratica habebunt indices per $4$ divisibiles sive formae $4n$; non-residua ea, quae sunt residua quadratica, habebunt indices formae $4n+2$; denique non-residuorum quadraticorum indices erunt partim formae $4n+1$, partim formae $4n+3$. Hoc modo classes quaternae quidem oriuntur, at distinctio inter binas posteriores non esset absoluta, sed ab electione radicis primitivae pro basi assumtae dependens; facile enim perspicitur, semissem radicum primitivarum non-residuo quadratico dato conciliare indicem formae $4n+1$, semissem alteram vero indicem formae $4n+3$. Quam ambiguitatem ut tollamus, supponemus semper talem radicem primitivam adoptari, pro qua index $\frac14(p-1)$ competat numero $+i$ (conf. art.~55, VI). Hoc pacto classificatio oritur, quam concinnius independenter a radicibus primitivis ita enunciare possumus.

Classis \emph{prima} contineat numeros $k$ eos, pro quibus fit $k^{\frac14(p-1)}\equiv1$; hi numeri sunt moduli residua biquadratica.

Classis \emph{secunda} contineat eos, pro quibus $k^{\frac14(p-1)}\equiv i$.

Classis \emph{tertia} eos, pro quibus $k^{\frac14(p-1)}\equiv-1$.

Classis \emph{quarta} denique eos, pro quibus $k^{\frac14(p-1)}\equiv-i$.

Classis tertia comprehendet non-residua biquadratica ea, quae sunt residua quadratica; inter secundam et quartam non-residua quadratica distributa erunt.

Numeris harum classium tribuemus resp. \emph{characteres biquadraticos} $0,1,2,3$. Si characterem $\lambda$ numeri $k$ secundum modulum $m$ ita definimus, ut sit exponens eius potestatis ipsius $i$, cui numerus $k^{\frac14(p-1)}$ congruus est, manifesto characteres secundum modulum $4$ congrui pro aequivalentibus habendi sunt. Ceterum haec notio tantisper ad modulos eos limitatur, qui sunt numeri primi: in continuatione harum disquisitionum ostendemus, quomodo etiam modulis compositis adaptari possit.

\Article{62}
Quo facilius inductio copiosa circa numerorum characteres adstrui possit, tabulam compendiosam hic adiungimus, cuius auxilio character cuiusvis numeri propositi respectu moduli, cuius norma valorem $157$ non transcendit, levi opera obtinetur, dummodo ad observationes sequentes attendatur.

\SourceLeaf{142}%
Quum character numeri compositi aequalis sit (sive secundum modulum $4$ congruus) aggregato characterum singulorum factorum, sufficit, si pro modulo dato characteres numerorum primorum assignare possumus. Porro quum characteres unitatum $-1,i,-i$ manifesto sint congrui numeris $\frac12(p-1),\frac14(p-1),\frac34(p-1)$ secundum modulum $4$, etiam sufficiet, characteres numerorum inter associatos primariorum exhibuisse. Denique quum moduli secundum modulum $m$ congrui eundem characterem habeant, sufficit, characteres talium numerorum in tabulam recipere, qui continentur in systemate residuorum absolute minimorum. Praeterea per ratiocinium simile ut in art.~58 demonstratur, si pro modulo $a+bi$ character numeri $A+Bi$ sit $\lambda$, pro modulo $a-bi$ autem $\lambda'$ sit character numeri $A-Bi$, semper esse $\lambda\equiv-\lambda'\pmod4$, sive $\lambda+\lambda'$ per $4$ divisibilem: quapropter sufficit, in tabulam recipere modulos, in quibus $b$ est vel $0$ vel positivus.

Ita e.~g. si quaeritur character numeri $11-6i$ respectu moduli $-5-6i$, substituimus loco horum numerorum hosce $11+6i,-5+6i$; dein determinamus (art.~43) residuum absolute minimum numeri $11+6i$ secundum modulum $-5+6i$, quod fit $-1-4i=-1\times(1+4i)$; quare quum pro modulo $-5+6i$ character ipsius $-1$ sit $30$, character numeri $1+4i$ autem, ex tabula, $2$, erit $32$ sive $0$ character numeri $11+6i$ pro modulo $-5+6i$, et proin per observationem ultimam etiam character numeri $11-6i$ pro modulo $-5-6i$. Perinde si quaeritur character numeri $-5+6i$ respectu moduli $11+6i$, illius residuum absolute minimum $1-5i$ resolvitur in factores $-i\;1+i,\;3-2i$, quibus respondent characteres $117,0,1$, unde character quaesitus erit $118$ sive $2$; idem character etiam numero $-5-6i$ respectu moduli $11-6i$ tribuendus est.

\[
\begin{array}{r|c|l}
\text{Modulus.}&\text{Character.}&\text{Numeri.}\\\hline
-3&3&1+i\\
+3+2i&3&1+i\\
+1+4i&1&-1+2i\\
&3&1+i\\
-5+2i&0&-1-2i\\
&1&1+i\\
&2&-1+2i\\
-1+6i&0&-3\\
&1&1+i,\ -1+2i
\end{array}
\]

\SourceLeaf{143}%
\[
\begin{array}{r|c|l}
\text{Modulus.}&\text{Character.}&\text{Numeri.}\\\hline
-1+6i&2&-1-2i\\
+5+4i&0&1+i\\
&1&-3\\
&3&-1+2i,\ -1-2i\\
-7&0&-3\\
&1&-1+2i,\ 3-2i\\
&2&1+i\\
&3&-1-2i,\ 3+2i\\
+7+2i&0&1+i,\ 3+2i,\ 3-2i,\ 1-4i\\
&1&-3\\
&2&-1-2i,\ 1+4i\\
&3&-1+2i\\
-5+6i&0&1+i,\ -3,\ 3+2i,\ 3-2i\\
&1&1-4i\\
&2&1+4i\\
&3&-1+2i,\ -1-2i\\
-3+8i&0&-1+2i,\ 3-2i,\ 1-4i\\
&1&1+i,\ 3+2i\\
&2&-3\\
&3&-1-2i,\ 1+4i,\ -5+2i\\
+5+8i&0&-1-2i\\
&1&-5-2i,\ -1+6i\\
&2&-1+2i,\ 3-2i\\
&3&1+i,\ -3,\ 3+2i,\ 1+4i,\ 1-4i\\
+9+4i&0&-1+2i,\ 3+2i\\
&1&1+i,\ -1-2i,\ 3-2i\\
&2&-3,\ 1+4i\\
&3&1-4i,\ -5+2i\\
-1+10i&0&1+i,\ -1+2i,\ -1-2i,\ 3+2i\\
&1&-3\\
&2&3-2i,\ -5+2i,\ 5-4i\\
&3&1+4i,\ 1-4i
\end{array}
\]

\SourceLeaf{144}%
\[
\begin{array}{r|c|l}
\text{Modulus.}&\text{Character.}&\text{Numeri.}\\\hline
+3+10i&1&1+i,\ -1-2i,\ 1-4i\\
&2&-3,\ 3+2i,\ 1+4i,\ -5-2i\\
&3&-1+2i,\ 3-2i\\
-7+8i&0&1+i,\ -7\\
&1&3+2i,\ 3-2i,\ 1-4i,\ -5-2i\\
&2&-1-2i,\ 1+4i,\ -5+2i,\ -1-6i\\
&3&-1+2i,\ -3,\ -1+6i\\
-11&0&-3\\
&1&1+i,\ 3-2i,\ 1+4i,\ -5+2i,\ 5+4i\\
&2&-1+2i,\ -1-2i\\
&3&3+2i,\ 1-4i,\ -5-2i,\ 5-4i\\
-11+4i&0&1+i,\ -1+2i,\ 3+2i,\ 5+4i\\
&1&-1-2i,\ -1+6i\\
&2&-5+2i\\
&3&-3,\ 3-2i,\ 1+4i,\ 1-4i,\ -5-2i\\
+7+10i&0&1+4i,\ 1-4i,\ -1+6i,\ -1-6i\\
&1&-1+2i,\ 3+2i,\ -5+2i\\
&2&1+i,\ 3-2i\\
&3&-1-2i,\ -3,\ -5-2i\\
+11+6i&0&1+i,\ -1+2i,\ -3,\ 1+4i,\ 1-4i,\ -7\\
&1&-1-2i,\ 3+2i,\ 3-2i\\
&2&-5-2i,\ -1+6i,\ 5-4i\\
&3&-5+2i,\ 5+4i,\ 7-2i.
\end{array}
\]

\Article{63}
Operam nunc dabimus, ut criteria communia modulorum, pro quibus numerus primus datus characterem eundem habet, per inductionem detegamus. Modulos semper supponimus primarios inter associatos, puta tales $a+bi$, pro quibus vel $a\equiv1,b\equiv0$, vel $a\equiv3,b\equiv2\pmod4$.

Respectu numeri $1+i$, a quo initium facimus, inductionis lex facilius arripitur, si modulos prioris generis (pro quibus $a\equiv1,b\equiv0$) a modulis posterioris generis (pro quibus $a\equiv3,b\equiv2$) separamus. Adiumento tabulae art. praec. invenimus respondere
\SourceLeaf{145}%
\[
\begin{array}{c|l}
\text{characterem}&\text{modulis primi generis.}\\\hline
0&5+4i,\ -7+8i,\ -7-8i,\ -11+4i\\
1&1-4i,\ -3+8i,\ -3-8i,\ 9+4i,\ -11\\
2&5-4i,\ -7,\ -11-4i\\
3&-3,\ 1+4i,\ 5+8i,\ 5-8i,\ 9-4i
\end{array}
\]
Si haec septemdecim exempla attente consideramus, in omnibus invenimus characterem $\equiv\frac14(a-b-1)\pmod4$.

Perinde respondet
\[
\begin{array}{c|l}
\text{character}&\text{modulis secundi generis.}\\\hline
0&3-2i,\ -1-6i,\ 7+2i,\ -5+6i,\ -1+10i,\ 11+6i\\
1&-5+2i,\ -1+6i,\ 7-2i,\ -1-10i,\ 3+10i\\
2&-1+2i,\ -5-2i,\ 3-10i,\ 7+10i\\
3&-1-2i,\ 3+2i,\ -5-6i,\ 7-10i,\ 11-6i
\end{array}
\]
In omnibus his viginti exemplis, levi attentione adhibita, invenitur character $\equiv\frac14(a-b-5)\pmod4$.

Facile has duas regulas in unam pro utroque modulorum genere valentem contrahere licet, si perpendimus, $\frac14bb$ esse pro modulis prioris generis $\equiv0$, pro modulis posterioris generis $\equiv1\pmod4$. Est itaque character numeri $1+i$ respectu moduli cuiusvis primi inter associatos primarii $\equiv\frac14(a-b-1-bb)\pmod4$.

Obiter hic annotare convenit, quum $(b+1)^2$ semper sit formae $8n+1$, sive $\frac14(2b+bb)$ par, characterem istum semper parem vel imparem fieri, prout $\frac14(a+b-1)$ par sit vel impar, quod quadrat cum regula pro charactere quadratico in art.~58 prolata.

Quum $\frac14(a-b-1),\frac14(a-b+3)$ sint integri, quorum alter par, alter impar, ipsorum productum par erit, sive $\frac18(a-b-1)(a-b+3)\equiv0\pmod4$. Hinc loco expressionis allatae pro charactere biquadratico haec quoque adoptari potest
\[
\frac14(a-b-1-bb)-\frac18(a-b-1)(a-b+3)=\frac18(-aa+2ab-3bb+1)
\]
quae forma eo quoque nomine se commendat, quod non restringitur ad modulos primarios, sed tantummodo supponit, $a$ esse imparem, $b$ parem: manifesto enim in hac suppositione vel $a+bi$, vel $-a-bi$ erit numerus inter associatos primarius, valorque istius formulae pro utroque modulo idem.

